\lecture{9}{Tue 17 Mar 2026 10:58}{I took notes today}

We want to decompose Lie algebras into more understandable components. The main tool we will use is the representation theory of Lie algebras, and we will largely consider the adjoint representation.

Let $G$ a Lie group. A representation of $G$ on a vector space $V$ is a Lie group homomorphism $\rho : G\to \Aut(V) $, where we note that $\Aut(V) $ is a Lie group: it's group structure is composition, and its manifold structure is given by viewing it as a vector space, choosing a basis, and using the isomorphism to $\mathbb{R} ^n$ to define a smooth structure. Given $g\in G$ and $v\in V$, we denote $\rho (g)(v)$ by $g\cdot v$.

For example, consider the action of $\operatorname{SO} (3)$ on $\mathbb{R} ^3$ by rotation. Since each $R\in\operatorname{SO} (3)$ is an automorphism of $\mathbb{R} ^3$, we have $\operatorname{SO} (3) \subseteq \Aut(\mathbb{R} ^3) $ is a Lie subgroup, and the inclusion is a Lie group homomorphism.

Now let $\mathfrak{g} $ be the Lie algebra of a Lie group $G$. Let $X\in \mathfrak{g} $, and let $\gamma (t)$ be a path in $G$ with $\gamma '(0)=X$. Let $G$ act on a vector space $V$, and let the action be given by $\rho : G \to \Aut(V) $. Then we define the action $\theta $ of the Lie algebra $\mathfrak{g} $ on the vector space $V$ to be
\[
	X\cdot v \coloneqq \left.\frac{\mathrm{d}}{\mathrm{d}t} \right|_{t=0} (\gamma (t)\cdot v)
.\]
This yields the commutative diagram
\[\begin{tikzcd}[row sep = 2.5em, column sep = 2.5em]
\mathfrak{g} \ar[" \theta  ", r] \ar[" \exp  "', d]  & \End(V) =\mathfrak{gl} (V)\ar[" \exp  ", d]  \\
G\ar[" \rho  "', r]  & \Aut(V) .
\end{tikzcd}\]
Note that this representation commutes with the Lie bracket. This motivates a more general definition: a representation of a Lie algebra $\mathfrak{g} $ is a Lie algebra homomorphism $\theta : \mathfrak{g}  \to \mathfrak{gl}_n(V)$. Note that the Lie algebra homomorphism condition implies that $\rho [X,Y]=[\rho (X),\rho (Y)]$.

We have a fairly canonical representation $\mathrm{ad} : \mathfrak{g}  \to \End(\mathfrak{g} ) $ given by the adjoint map:
\[
	\mathrm{ad} : X \mapsto \mathrm{ad}_X = [X,-]
.\]
Recall that Jacobi's identity implies $\mathrm{ad}_{[X,Y]} =[\mathrm{ad} _X,\mathrm{ad} _Y]$, hence this is a Lie algebra homomorphism. We call this the \emph{adjoint representation} of $\mathfrak{g} $.

We would like to regard $\mathfrak{g} $ as a representation of itself by breaking it into smaller pieces.

A vector space $V$ together with a representation of $\mathfrak{g} $ on $V$ is called a \emph{$\mathfrak{g} $-module}. Explicitly, this means $[X,Y]v=X(Yv)-Y(Xv)$. We will generally represent a representation by the name of the vector space. This allows us to consider a direct sum of representations: given representations $V,W$ of $\mathfrak{g}$, define the representation of $\mathfrak{g} $ on $V\oplus W$ as
\[
	X\cdot (v,w) = (X\cdot v,X\cdot w)
.\]
Given a subspace $W\subseteq V$ of a representation $V$ of $\mathfrak{g} $ together with the property that $\mathfrak{g} W \subseteq W$, define the \emph{subrepresentation} of $\mathfrak{g} $ on $W$ to be the representation of $\mathfrak{g} $ on $V$ but restricted to $W$. Given a subrepresentation $W$ of $V$, define the quotient representation $V/W$ to be $X\cdot [v]\coloneqq [X\cdot v]$. This is well-defined. Given a representation $V$, define the \emph{dual representation} $\rho ^*$ on the dual $V^*$ to be
\[
	\rho^*(X) \coloneqq -\rho (X)^t
\]
for $-^t$ the transpose.

Consider the adjoint representation of $\mathfrak{g} $ on $\mathfrak{g} $. Suppose we have a subrepresentation $\mathfrak{i} $ of this. We call this an \emph{ideal}, and it is precisely a Lie subalgebra $\mathfrak{i} \subseteq \mathfrak{g} $ such that $[\mathfrak{g} ,\mathfrak{i} ]\subseteq \mathfrak{i} $.

For example, consider the ideal $[\mathfrak{g} ,\mathfrak{g} ]$. This is immediately an ideal, and the quotient $\mathfrak{g} ^{ab} \coloneqq \mathfrak{g} /[\mathfrak{g} ,\mathfrak{g} ]$ is abelian.

This motivates the \emph{derived series}: define $\mathfrak{g} ^{(n)} \coloneqq [\mathfrak{g} ^{(n-1)} ,\mathfrak{g} ^{(n-1)} ]$. Then $\mathfrak{g} ^{(n-1)} \subseteq \mathfrak{g} ^{(n)} $, and this series is the \emph{derived series}. Each $\mathfrak{g} ^{(n)} $ is an ideal, and the quotient by that ideal is an abelian Lie algebra. The series stabilizes. Apparently $\mathfrak{g} = \mathfrak{g} ^{ab} \ltimes [\mathfrak{g},\mathfrak{g} ]$.

We want to write all Lie algebras as $\mathfrak{g} = \mathfrak{s} \ltimes \mathfrak{r} $, where $\mathfrak{s} $ is \emph{semi-simple}, meaning it has no nonzero Abelian ideal, and $\mathfrak{r} $ is a \emph{solvable radical}, meaning it's a solvable ideal (its derived series converges to 0) and it is the maximum solvable ideal. Classifying Lie algebras thus apparently amounts to classifying semi-simple Lie algebras. For example, the ???

look up Engel's theorem for nilpo;tent representation also something about $\mathfrak{gl} _{n-k} $ being $\mathfrak{sl} _{n-k} $ as semisimple and upper triangular as solvable radical.
